\documentclass[11pt]{ctexart}
\usepackage[a4paper,margin=1in]{geometry}
\usepackage{graphicx}
\usepackage{xcolor}
\usepackage{hyperref}
\definecolor{refgray}{gray}{0.45}
\title{\textbf{Market Views｜全球投行叙事汇编}\\[0.4em]{\large 通胀降温与地缘扰动交织，AI资本开支加速但回报验证成新焦点}}
\author{KC桌面 自动生成}
\date{260721}
\begin{document}
\maketitle
\tableofcontents
\newpage
\section{一页摘要}
\begin{itemize}
  \item 美国核心通胀趋势性回落，美联储7月加息概率基本消失，但中东局势升级推升能源成本，对亚洲进口国形成新的外部约束。
  \item 全球超大规模云厂商AI资本开支持续上行，2026年底预计接近8700亿美元，同比增长约77\%，但市场关注点从规模转向回报率。
  \item 中国电池消费税重启（2\%→4\%）加速行业洗牌，头部企业利润影响仅1-6\%，二线厂商面临亏损风险，行业集中度提升。
  \item 铜市场短期实物收紧（中国库存降至16.7万吨、LME注销仓单上升），但宏观不确定性压制价格上行空间，AI叙事正成为定价新变量。
  \item 中国消费市场呈现收入改善但支出谨慎的背离，居民可支配收入增速回升至5.6\%，而社零增速降至0.2\%，储蓄率微升至32.6\%。
  \item 茅台一年内两次调价，出厂价累计上调17\%，显示公司正主动管理渠道利润和品牌定价能力，但批发价反应温和。
  \item 韩国市场上涨带来巨额名义财富，但高集中度与高杠杆制约消费转化，实际拉动效果存疑。
  \item 印度估值已重置至疫情低位，但盈利下调周期是否见底、外资流入能否持续，机构看法不一。
\end{itemize}
\section{全球投行叙事汇编}
今日纳入 31 篇投行/券商研究，按 4 个市场、资产与产业主题系统整理。
\newpage
\subsection{宏观与策略：通胀降温与地缘扰动下的资产定价再校准}
\textbf{全球宏观主线正从通胀恐慌转向增长分化与地缘风险再定价。美国核心通胀趋势性回落为美联储政策转向提供空间，但中东局势升级推升能源成本，对亚洲进口国形成新的外部约束。中国内需疲弱与韩国财富效应传导不畅凸显增长结构分化，而印度估值重置与AI板块盈利支撑则为资金轮动提供结构性机会。}
\subsubsection*{主流共识}
\begin{itemize}
  \item 全球除美国外核心通胀已接近目标，美国核心PCE高估真实通胀压力，美联储7月加息概率基本消失。
  \item 中东地缘不确定性推升能源价格，但供需两端均存在快速逆转可能，油价基准判断仍是中期回落。
  \item AI相关板块调整更多是动量退潮与拥挤度释放，而非基本面逆转，半导体盈利支撑依然强劲。
\end{itemize}
\subsubsection*{分歧与条件差异}
\begin{itemize}
  \item 中国增长节奏：二季度GDP低于目标下限，但市场对三季度财政加速的力度与效果存在分歧。
  \item 韩国财富效应：市场上涨带来巨额名义财富，但高集中度与高杠杆制约消费转化，实际拉动效果存疑。
  \item 印度市场：估值已重置至疫情低位，但盈利下调周期是否见底、外资流入能否持续，机构看法不一。
  \item 人民币定价：NOM模型显示中间价预测下调，但逆周期因子调整方向与误差收敛是否预示趋势转变，仍需观察。
\end{itemize}
\subsubsection*{机构观点}
\paragraph{Bernstein} 印度市场估值已重置至18.6倍远期PE，低于十年均值0.7个标准差，金融与科技板块PE回到疫情低位。盈利下调周期尚未见底，但大盘股及材料、消费板块出现早期复苏信号。7月外资净流入18亿美元，全球新兴市场基金对印度低配从-2.7\%收窄至-2.1\%。维持中性，但机会组合正在变宽。
{\small\color{refgray}数据：MSCI印度远期PE 18.6倍，低于十年均值0.7个标准差\par}
{\small\color{refgray}数据：金融与科技板块PE低于均值1.9和1.5个标准差\par}
{\small\color{refgray}数据：7月外资净流入18亿美元\par}
{\small\color{refgray}数据：全球新兴市场基金对印度低配从-2.7\%收窄至-2.1\%\par}
{\small\color{refgray}边际变化：印度估值从23.2倍PE高点回落至18.6倍，外资流出趋势逆转，大盘股盈利修正出现早期改善。\par}
\paragraph{GS} 韩国市场上涨为家庭带来相当于GDP约17\%的财富增值，但每100韩元财富仅带动1.6韩元额外消费，全年GDP拉动约0.3个百分点。前20\%高收入家庭持有64\%股票资产，50岁以上群体持有73\%，高储蓄率（37\%）与高杠杆（偿债负担占收入27\%）制约消费转化。央行加息可能完全抵消财富效应。
{\small\color{refgray}数据：家庭财富增值相当于GDP的17\%\par}
{\small\color{refgray}数据：每100韩元财富带动1.6韩元额外消费\par}
{\small\color{refgray}数据：前20\%高收入家庭持有64\%股票资产\par}
{\small\color{refgray}数据：50岁以上群体持有73\%股票资产\par}
{\small\color{refgray}边际变化：市场上涨幅度创历史新高，但财富效应传导系数低于预期，加息周期启动进一步压制消费转化。\par}
\paragraph{GS} 中国二季度实际GDP同比4.3\%，低于市场预期4.5\%及政府目标下限。6月工业增加值同比5.3\%，但零售仅增1.0\%，固定资产投资累计同比-5.7\%，生产与需求裂口扩大。贸易顺差创纪录1256亿美元，但半导体进口量仅增6.6\%，出口量微降0.5\%，量价拆分显示含金量有限。预计7月政治局会议将加速财政支出，8000亿元政策性金融工具可拉动GDP约0.5个百分点。
{\small\color{refgray}数据：二季度实际GDP同比4.3\%，低于预期4.5\%\par}
{\small\color{refgray}数据：6月工业增加值同比5.3\%，零售同比1.0\%\par}
{\small\color{refgray}数据：固定资产投资累计同比-5.7\%\par}
{\small\color{refgray}数据：6月贸易顺差1256亿美元创纪录\par}
{\small\color{refgray}边际变化：GDP增速从一季度的5.0\%降至4.3\%，低于政府目标下限，触发三季度财政加速预期。\par}
\paragraph{GS} 中东局势升级推升布伦特原油突破此前预测路径，但油价不确定性双向存在。若局势缓和，海湾出口可快速恢复；若基础设施受攻击，价格可能重回100美元以上。基准判断仍是中期回落。美国核心PCE约3.3\%存在测量偏差，G10除美国外核心通胀仅2.1\%。6月非农低于预期，薪资跟踪指标已降至与通胀目标匹配水平。预计下半年美国消费将放缓。
{\small\color{refgray}数据：布伦特原油突破GS此前预测的2026Q4 80美元/桶\par}
{\small\color{refgray}数据：G10除美国外核心通胀仅2.1\%\par}
{\small\color{refgray}数据：美国核心PCE约3.3\%\par}
{\small\color{refgray}数据：6月非农低于预期\par}
{\small\color{refgray}边际变化：地缘风险推升油价但通胀趋势性降温，美联储7月加息概率基本消失，市场主线从通胀恐慌转向增长分化。\par}
\paragraph{JPM} AI相关板块经历剧烈调整，KOSPI从高点跌25\%，费城半导体跌20\%，但MSCI全球指数仅较高点回落1\%-2\%。半导体相对价格与盈利出现罕见裂口，盈利支撑依然强劲，产能增量要到2028年后才到来。RSI逼近超卖区域，动量因子拥挤度显著下降。通胀方面，美国CPI三个月年化从5月8.2\%降至6月2.8\%，油价同比下跌约25\%的基数效应将继续压制通胀。
{\small\color{refgray}数据：KOSPI从高点下跌25\%\par}
{\small\color{refgray}数据：费城半导体指数下跌20\%\par}
{\small\color{refgray}数据：MSCI全球指数仅较高点回落1\%-2\%\par}
{\small\color{refgray}数据：美国CPI三个月年化从8.2\%降至2.8\%\par}
{\small\color{refgray}边际变化：AI动量退潮进入成熟阶段，但盈利支撑与估值吸引力上升，半导体板块可能很快找到买盘。\par}
\paragraph{NOM} TACO不确定性指标从7月6日接近0标准差升至2.0标准差，主因美伊停火备忘录失效后霍尔木兹海峡事实性关闭。油价和天然气需再涨10\%才会触及3.2标准差阈值，当前反映不确定性积累而非临界点。美元/人民币中间价模型预测值下调110基点至6.7824，计入逆周期因子后为6.7862，较实际中间价低72基点，误差收敛表明定价预期正在重新校准。
{\small\color{refgray}数据：TACO指标从接近0标准差升至2.0标准差\par}
{\small\color{refgray}数据：油价和天然气需再涨10\%才触及3.2标准差\par}
{\small\color{refgray}数据：美元/人民币中间价模型预测值6.7824，下调110基点\par}
{\small\color{refgray}数据：计入逆周期因子后预测6.7862，较实际低72基点\par}
{\small\color{refgray}边际变化：美伊停火备忘录五周内失效，地缘不确定性快速攀升但尚未触发紧急机制；人民币中间价模型误差收敛，定价预期趋于温和。\par}
\subsubsection*{关键数据}
\begin{itemize}
  \item MSCI印度远期PE 18.6倍，低于十年均值0.7个标准差
  \item 韩国家庭财富增值相当于GDP的17\%，每100韩元仅带动1.6韩元消费
  \item 中国二季度GDP同比4.3\%，低于政府目标下限
  \item 6月中国贸易顺差1256亿美元创纪录
  \item G10除美国外核心通胀仅2.1\%，美国核心PCE 3.3\%
  \item 美国CPI三个月年化从8.2\%降至2.8\%
  \item KOSPI从高点下跌25\%，费城半导体跌20\%
  \item NOM TACO指标从0升至2.0标准差
  \item 美元/人民币中间价模型预测下调110基点
\end{itemize}
\begin{figure}[htbp]
\centering
\includegraphics[width=0.88\linewidth]{figures/fig_001.jpg}
\caption{Exhibit 1 - Exhibit 1-Exhibit 3) EXHIBIT 1: India valuations are now cheap both on absolute and relative basis, which trade at 18.6x on fwd PE and at 3.3x on PB. Relative to EM, India's PB has fallen near COVID lows India relative to EM-12m}
\end{figure}
\begin{figure}[htbp]
\centering
\includegraphics[width=0.88\linewidth]{figures/fig_027.jpg}
\caption{Exhibit 1 - Exhibit 2: Korean households' equity assets could have risen by \$15\textbackslash{}\%\$ of GDP from the end of last year Percent of GDP}
\end{figure}
\begin{figure}[htbp]
\centering
\includegraphics[width=0.88\linewidth]{figures/fig_057.jpg}
\caption{Exhibit 1 - Exhibit 2: We Still Expect US Consumer Spending to Slow in H2 \textbackslash{}*Both series adjusted for AI-related mismeasurement. \textbackslash{}*\textbackslash{}*Real DPI with taxes converted from accrual to cash.}
\end{figure}
\begin{figure}[htbp]
\centering
\includegraphics[width=0.88\linewidth]{figures/fig_062.jpg}
\caption{Figure 1 - Figure 2: SOX Index ytd}
\end{figure}
{\small\color{refgray}本节综合 7 篇研究；涉及 Bernstein、GS、JPM、NOM。}
\newpage
\subsection{AI基础设施：资本开支加速与结构性瓶颈}
\textbf{AI基础设施投资正从规模扩张转向效率与回报验证，中国本土芯片与全球存储器成为关键瓶颈，而人形机器人模型架构的迭代预示下一阶段竞争焦点。}
\subsubsection*{主流共识}
\begin{itemize}
  \item 全球超大规模云厂商AI资本开支持续上行，2026年底预计接近8700亿美元，同比增长约77\%。
  \item 中国主要云厂商AI基础设施投入被大幅上修，2026-2028年预测分别上调37\%、44\%和55\%，驱动本土芯片需求。
  \item AI应用渗透加速，token消耗量稳步提升，但货币化路径存在时间差，市场关注点从资本开支规模转向回报率。
\end{itemize}
\subsubsection*{分歧与条件差异}
\begin{itemize}
  \item 存储器周期：MS认为短缺加剧且将持续数年，去规格化拉长周期而非终结；部分市场观点担忧长期协议压低价格峰值。
  \item 中国AI芯片替代节奏：GS认为本土供应商份额提升刚刚开始，非本土供应商仍占超50\%份额；但市场对技术路线切换和价格竞争存在担忧。
  \item 人形机器人技术路径：Bernstein强调从VLA向世界动作模型（WAM）的范式转变，认为模型架构和数据整合能力将取代数据垄断成为新护城河；部分观点仍聚焦于硬件进展。
\end{itemize}
\subsubsection*{机构观点}
\paragraph{Bernstein} 人形机器人行业正从VLA模型转向世界动作模型（WAM），后者通过预测物理可行结果而非直接输出动作，实现更优的自主规划与任务泛化。当前技术前沿在运动控制、操作、自主性、大脑模型和数据模态五个维度均有突破，但长序列自主规划和技能迁移仍是下一阶段战场。WAM强调数据多样性和模型架构，将削弱数据垄断壁垒。
{\small\color{refgray}数据：Unitree机器人能在波兰森林追逐野猪\par}
{\small\color{refgray}数据：Boston Dynamics Atlas做出欺骗性足球射门\par}
{\small\color{refgray}数据：Physical Intelligence π0.7模型引入世界模型后复杂任务表现提升\par}
{\small\color{refgray}边际变化：行业焦点从'能不能动'转向'能不能想'，模型架构竞争取代数据规模竞争成为新护城河。\par}
\paragraph{Bernstein} 中国互联网板块估值回调至2022-2023年低谷，但AI推理成本担忧过度。腾讯Hy3模型发布、阿里云增长符合预期、电商利润好于担忧等信号触发近期反弹。AI推理成本仅在代理交易量和GMV起飞时才会指数级增长，从用户参与到货币化存在时间差。阿里云6月季度收入增速加速至45\%左右，下半年计算价格将继续上调。
{\small\color{refgray}数据：覆盖公司年初至今平均下跌约16\%\par}
{\small\color{refgray}数据：腾讯估值倍数约11-12倍\par}
{\small\color{refgray}数据：阿里云6月季度收入增速约45\%\par}
{\small\color{refgray}数据：阿里云增量现金利润率35-40\%\par}
{\small\color{refgray}边际变化：市场从担忧AI资本开支回报转向关注货币化时间差，估值已充分反映坏消息。\par}
\paragraph{GS} 美国线上杂货渗透率仍处低位，预计2025-2030年行业复合增速达11\%，2030年市场规模接近4000亿美元。当日配送为增长最快赛道（+16\% CAGR），驱动力来自现有用户购买频率提升而非新增用户。近50\%消费者过去一年增加快速配送订单次数，配送时效从小时级向分钟级跨越将解锁更高频次。
{\small\color{refgray}数据：2025-2030年行业复合增速+11\%\par}
{\small\color{refgray}数据：2030年市场规模接近4000亿美元\par}
{\small\color{refgray}数据：当日配送2025-2030年复合增速+16\%\par}
{\small\color{refgray}数据：近50\%消费者增加快速配送订单次数\par}
{\small\color{refgray}边际变化：行业增长驱动力从'拉新'转向'复购'，配送时效缩短和自动化履约中心规模化是频率提升的物理前提。\par}
\paragraph{GS} 全球夏季出行需求稳定健康，6月预订趋势超季初预期，但呈现K型消费分化：高端市场定价上移，大众消费者削减其他开支挤出旅行预算。约三分之一受访者减少外出就餐、线上购物和流媒体订阅以腾出旅行预算。AI在旅行规划中使用率显著提升，1H26美国超56\%休闲旅客使用AI规划旅行（去年同期33\%），但通过AI完成预订比例仍低。
{\small\color{refgray}数据：1H26美国超56\%休闲旅客使用AI规划旅行（去年同期33\%）\par}
{\small\color{refgray}数据：约三分之一受访者削减其他开支为旅行腾预算\par}
{\small\color{refgray}数据：单次旅行支出上升但总预订人数可能同比下降\par}
{\small\color{refgray}数据：北美主办城市赛事期间短租均价同比上涨60\%以上\par}
{\small\color{refgray}边际变化：消费者预算再分配而非消费升级，AI在旅行规划中渗透率翻倍但预订转化仍低。\par}
\paragraph{GS} 中国数据中心电气公司股价回调约29\%，核心担忧是技术路线切换和国内价格竞争，但海外扩张和800VDC产品商业化是触发重新定价的关键变量。Kstar大功率UPS产线接近满负荷，已获欧洲数据中心超2亿元订单（最终租给美国超大规模客户），证明中国供应商进入全球供应链核心环节的能力。800VDC架构最早2027下半年规模落地。
{\small\color{refgray}数据：Kstar、Kehua、Megmeet股价平均下跌约29\%\par}
{\small\color{refgray}数据：Kstar大功率UPS产线接近满负荷运转\par}
{\small\color{refgray}数据：Kstar获欧洲数据中心超2亿元UPS订单\par}
{\small\color{refgray}数据：800VDC架构最早2027下半年规模落地\par}
{\small\color{refgray}边际变化：海外订单落地速度决定估值修复起点，中国供应商正从认证阶段进入核心供应链。\par}
\paragraph{GS} 一家中国光模块公司2Q26净利润指引中值较预期高出44\%，核心驱动力是800G和1.6T光模块出货加速，而非整体需求突然放大。光学芯片供应瓶颈正在系统性突破，产能扩张配合出货节奏。2026-2028年盈利预测分别上调7\%、2\%和2\%，增长逻辑将从规模扩张转向规模互联（scale-out到scale-up/scale-across）。
{\small\color{refgray}数据：2Q26净利润指引中值较预期高出44\%\par}
{\small\color{refgray}数据：2Q26净利润同比增长78\%至163\%，环比增长68\%\par}
{\small\color{refgray}数据：2026年营收预期上调6\%至514.6亿\par}
{\small\color{refgray}数据：毛利率从2023年31\%提升至2025年47.8\%\par}
{\small\color{refgray}边际变化：光学芯片供应改善驱动800G/1.6T加速放量，产品结构升级而非总量扩张是超预期主因。\par}
\paragraph{GS} 中国AI基础设施投入加速上行，本土芯片供应商份额提升刚刚开始。一家本土GPU芯片设计公司收入预计2025-2030年以108\%年复合增速增长，芯片出货量从2025年4.5万片攀升至2030年超100万片。全球领先AI芯片供应商在中国仍占超50\%市场份额，本土替代空间显著。中国云厂商AI资本开支2026-2028年预测分别上调37\%、44\%和55\%。
{\small\color{refgray}数据：2026-2028年中国云厂商AI资本开支预测分别上调37\%、44\%、55\%\par}
{\small\color{refgray}数据：本土GPU芯片公司收入2025-2030年CAGR 108\%\par}
{\small\color{refgray}数据：芯片出货量从2025年4.5万片增至2030年超100万片\par}
{\small\color{refgray}数据：非本土AI芯片供应商在中国仍占超50\%市场份额\par}
{\small\color{refgray}边际变化：中国AI应用渗透加速推动基础设施投资，本土替代从'刚刚开始'进入加速阶段。\par}
\paragraph{JPM} AI基础设施叙事未变，但市场关注点从资本开支规模转向货币化回报。预计到2026年底AI相关资本开支接近8700亿美元（同比增长约77\%），超大规模厂商贡献约7500亿。JPM对2027年预测更为积极：谷歌资本开支预计增长54\%至约3000亿美元，亚马逊和Meta分别增长42\%。动量因子从极端拥挤水平回落，但波动可能持续。
{\small\color{refgray}数据：2026年底AI相关资本开支预计接近8700亿美元\par}
{\small\color{refgray}数据：同比增长约77\%\par}
{\small\color{refgray}数据：超大规模厂商贡献约7500亿美元\par}
{\small\color{refgray}数据：谷歌2027年资本开支预计增长54\%至约3000亿美元\par}
{\small\color{refgray}边际变化：市场从'要不要投'进入'投了能赚回多少'阶段，货币化信号若浮现将支撑盈利和现金流预期。\par}
\paragraph{MS} 存储器已从AI需求受益者变成AI供给瓶颈。数据中心内存价格三季度环比上涨至少25\%，高于预期，短缺加剧。AI消耗DRAM过多，PC和智能手机出货受抑制。长期协议和去规格化改变周期形状而非终结方向，提供更长盈利可见度。客户愿为六周加急订单支付溢价，反映真实产能挤兑而非囤货。
{\small\color{refgray}数据：数据中心内存价格三季度环比上涨至少25\%\par}
{\small\color{refgray}数据：AI相关资本开支到2026年底接近8700亿美元\par}
{\small\color{refgray}数据：客户愿为六周加急订单支付高于二季度合同价\par}
{\small\color{refgray}数据：英伟达削减机架中LPDDR5内存含量以优化使用\par}
{\small\color{refgray}边际变化：存储器从AI需求受益者变为供给瓶颈，短缺加剧且将持续数年，去规格化确认短缺长期性。\par}
\subsubsection*{关键数据}
\begin{itemize}
  \item 2026年底AI相关资本开支预计接近8700亿美元，同比增长约77\%
  \item 中国云厂商AI资本开支2026-2028年预测分别上调37\%、44\%、55\%
  \item 数据中心内存价格三季度环比上涨至少25\%
  \item 本土GPU芯片出货量从2025年4.5万片增至2030年超100万片
  \item 800G/1.6T光模块公司2Q26净利润指引中值较预期高出44\%
  \item 美国线上杂货2025-2030年复合增速+11\%，2030年规模近4000亿美元
  \item 1H26美国超56\%休闲旅客使用AI规划旅行（去年同期33\%）
  \item 非本土AI芯片供应商在中国仍占超50\%市场份额
\end{itemize}
\begin{figure}[htbp]
\centering
\includegraphics[width=0.88\linewidth]{figures/fig_072.jpg}
\caption{Figure 2 - Figure 2: AI Capex (LTM Actuals vs Estimates) Figure 3: Hyperscaler Income Statement}
\end{figure}
\begin{figure}[htbp]
\centering
\includegraphics[width=0.88\linewidth]{figures/fig_076.jpg}
\caption{Exhibit 1 - Exhibit 1: DRAM multiples have corrected sharply, spot pricing is continuing to move higher Aren't the long-term agreements capping upside? Our view is that the long-term agreements are important to the stocks, but more because t}
\end{figure}
\begin{figure}[htbp]
\centering
\includegraphics[width=0.88\linewidth]{figures/fig_051.jpg}
\caption{Exhibit 1 - Exhibit 2: Earnings revision}
\end{figure}
\begin{figure}[htbp]
\centering
\includegraphics[width=0.88\linewidth]{figures/fig_017.jpg}
\caption{Exhibit 2 - Exhibit 2: We expect the industry to reach \textbackslash{}\textasciitilde{}\textbackslash{}\$400bn in sales by 2030... US grocery spend (\textbackslash{}\$bn) \#\# Across the three key categories of online grocery spend, we forecast:}
\end{figure}
{\small\color{refgray}本节综合 9 篇研究；涉及 Bernstein、GS、JPM、MS。}
\newpage
\subsection{能源与材料：电池税改重塑竞争格局，铜与电气化叙事分化}
\textbf{中国电池消费税重启加速行业洗牌，头部企业凭借利润垫与海外布局展现韧性；铜实物基本面收紧但宏观情绪退潮，AI叙事正成为定价新变量；欧洲电气化目标明确但电网投资缺口巨大，储能与电网设备需求进入加速期。}
\subsubsection*{主流共识}
\begin{itemize}
  \item 中国电池消费税（2\%→4\%）对终端需求影响有限，但将加速二线厂商出清，行业集中度提升。
  \item 铜市场短期实物收紧（中国库存降至16.7万吨、LME注销仓单上升），但宏观不确定性（中东局势、美联储政策）压制价格上行空间。
  \item 欧洲电气化目标（2040年电力占比46\%）与AI数据中心需求将推动电力投资周期，电网老化与需求爆发的时间差是核心矛盾。
\end{itemize}
\subsubsection*{分歧与条件差异}
\begin{itemize}
  \item 电池消费税影响程度：Bernstein认为成本增幅1-3\%不足以改变需求曲线，GS强调CATL等头部企业利润影响仅1-6\%，而二线厂商可能面临亏损风险。
  \item 铜价驱动逻辑：GS模型显示AI叙事解释力上升，但数据中心仅占铜需求1\%，市场对AI预期是否过度定价存在分歧。
  \item 美国关税对供应链的影响：GS高频数据显示洛杉矶港到港量短期反弹后将继续下行，但市场对“抢运”与“观望”交替节奏的持续性看法不一。
  \item 瑞萨存储器接口调查影响：Bernstein认为最差情景下对营收影响仅0.8\%，但市场担忧调查可能改变行业定价结构。
\end{itemize}
\subsubsection*{机构观点}
\paragraph{Bernstein} 储能电池正从电动汽车产能附属角色独立，三星SDI为特斯拉建专用储能产线、LG供应谷歌最大光储项目标志定制化供应关系形成。中国电池消费税（2\%→4\%）对终端价格影响仅1-3\%，但豁免钠离子/固态电池至2028年底，政策意图是引导产能向下一代技术转移。瑞萨存储器接口调查影响可控，即使毛利率从65\%降至55\%，对整体营收影响仅0.8\%。
{\small\color{refgray}数据：三星SDI为特斯拉建专用储能产线，200Ah LFP电芯2026年底量产，300Ah 2027年Q2投产\par}
{\small\color{refgray}数据：中国LFP电池组价格79.3美元/kWh，NMC 101.5美元/kWh\par}
{\small\color{refgray}数据：4\%消费税下电池组成本上升约3\%，储能系统成本上升1-2\%\par}
{\small\color{refgray}数据：瑞萨存储器接口业务占2025年营收约6\%，毛利率65\%\par}
{\small\color{refgray}边际变化：储能电池从借用动力电池产能转向自建专用制造体系，消费税政策从普惠免税转向技术差异化征税，瑞萨调查市场担忧被量化分析证伪。\par}
\paragraph{GS} 铜实物基本面收紧但宏观情绪退潮，AI叙事正重塑定价逻辑——AI因子对铜价解释力持续上升，尽管数据中心仅占全球铜需求1\%。中国电池消费税下CATL最具韧性（单位净利润109元/kWh，海外收入占比超30\%），在4\%税率完全自担下利润影响仅2-13\%。欧洲电气化计划支撑超采纳情景，未来十年需2.2-3.5万亿欧元投资，电力需求2029年起年均增长4-5\%。美国关税影响追踪显示洛杉矶港到港量短期反弹后将继续下行，中国出发货量同比降25.5\%。
{\small\color{refgray}数据：中国铜库存降至16.7万吨，处于季节性低位\par}
{\small\color{refgray}数据：LME铜价波动区间13,400-13,640美元/吨\par}
{\small\color{refgray}数据：CATL 2025年单位净利润109元/kWh，4\%税率对应税额11-24元/kWh\par}
{\small\color{refgray}数据：欧洲2030年储能装机目标200GW，2040年500GW\par}
{\small\color{refgray}边际变化：铜定价从中国需求+美元传统框架转向AI叙事+实物库存双轨，电池消费税从行业普惠转向龙头优势放大，欧洲电气化从政策目标转向具体KPI与投资测算。\par}
\paragraph{NOM} 中国锂电池消费税政策正常化，每辆电动车成本增加约400-500元，对整车制造影响有限。头部电池企业可通过成本转嫁消化冲击，二线厂商议价能力弱面临更大利润压力。新技术路线（钠离子、固态电池）获免税至2028年底，出口导向型厂商相对受益。
{\small\color{refgray}数据：2\%消费税对应每kWh成本增加约7元\par}
{\small\color{refgray}数据：60kWh电动车单车成本增加约400-500元\par}
{\small\color{refgray}数据：头部企业海外收入占比超30\%，二线厂商国内销售占比超98\%\par}
{\small\color{refgray}边际变化：电池行业从十年免税期进入税收正常化阶段，政策从普惠转向技术差异化，二线厂商利润空间被进一步压缩。\par}
\subsubsection*{关键数据}
\begin{itemize}
  \item 中国LFP电池组价格79.3美元/kWh，NMC 101.5美元/kWh
  \item CATL 2025年单位净利润109元/kWh，二线厂商仅8元/kWh
  \item 中国铜库存16.7万吨，处于季节性低位
  \item 欧洲2030年储能装机目标200GW，2040年500GW
  \item 未来十年欧洲电气化投资需2.2-3.5万亿欧元
  \item 中国至美国集装箱船数量同比降25.5\%
  \item 三星SDI为特斯拉建专用储能产线，价值3-5万亿韩元
  \item 瑞萨存储器接口业务毛利率65\%，最差情景营收影响0.8\%
\end{itemize}
\begin{figure}[htbp]
\centering
\includegraphics[width=0.88\linewidth]{figures/fig_050.jpg}
\caption{Exhibit 1 - Exhibit 1: A 4\% tax rate would imply an estimated unit tax of Rmb11–24/kWh for our covered companies, compared with their 2025 unit net profit of Rmb8–109/kWh.}
\end{figure}
\begin{figure}[htbp]
\centering
\includegraphics[width=0.88\linewidth]{figures/fig_037.jpg}
\caption{Exhibit 1 - Exhibit 3: China Copper Inventories Have Fallen to Seasonal Lows}
\end{figure}
\begin{figure}[htbp]
\centering
\includegraphics[width=0.88\linewidth]{figures/fig_016.jpg}
\caption{EXHIBIT 1 - EXHIBIT 1: Scenario Analysis on Renesas' Memory interface Business}
\end{figure}
\begin{figure}[htbp]
\centering
\includegraphics[width=0.88\linewidth]{figures/fig_067.jpg}
\caption{Exhibit 3 - Exhibit 4: TEU growth was decelerated in the most recent week Daily TEU from China to the US, YoY Data as of 7/16/2026. Represents the aggregated container volume, measured in twenty-foot equivalent units (TEU), of vessels depart}
\end{figure}
{\small\color{refgray}本节综合 8 篇研究；涉及 Bernstein、GS、NOM。}
\newpage
\subsection{消费与区域市场：结构分化中的主动调整}
\textbf{中国消费市场呈现收入改善但支出谨慎的背离，居民可支配收入增速回升至5.6\%，而社零增速降至0.2\%，储蓄率微升至32.6\%。在此背景下，消费品牌和区域银行均展现出主动调整资产结构或产品策略的迹象，而非被动等待宏观改善。}
\subsubsection*{主流共识}
\begin{itemize}
  \item 中国消费动能环比放缓，居民储蓄意愿仍强，超额存款达61万亿元。
  \item 8-10元价格带是中国啤酒高端化的核心战场，大单品策略在总量见顶市场中仍有效。
  \item 茅台一年内两次调价，出厂价累计上调17\%，显示公司正主动管理渠道利润和品牌定价能力。
  \item 安踏多品牌策略在行业减速中展现韧性，二季度核心品牌和FILA均实现低个位数增长。
\end{itemize}
\subsubsection*{分歧与条件差异}
\begin{itemize}
  \item 关于消费复苏路径：GS认为工资增长正在温和修复（二季度同比4.7\%），但市场感知与数据存在温差；而NOM关注金价波动对高端珠宝销售的短期扰动，认为稳定后情绪可能改善。
  \item 关于区域银行策略：GS强调TFC等银行主动压缩低收益消费贷款、转向商业贷款和手续费收入，而非等待利率环境改善；这与市场普遍预期的被动等待形成对比。
  \item 关于茅台调价信号：GS认为一年两涨打破了旧逻辑，定价权从市场供需转向公司战略主导；但市场对调价后的批发价反应温和（原箱飞天从1650元升至1720元），显示渠道消化仍需时间。
\end{itemize}
\subsubsection*{机构观点}
\paragraph{GS} 美国区域银行二季度财报显示，TFC、FITB、RF三家银行核心PPNR均略超预期，但市场对指引细节更敏感。TFC下调2026年净利息收入预期至持平/微降，同时将手续费收入增速上调至10\%，并主动压缩低收益间接汽车贷款、退出船舶和房车贷款，将资本转向商业贷款和国际业务。2026年ROTCE预期从14\%上调至14\%以上，长期目标16-18\%。当前估值约1.6倍有形市净率，低于同业1.8倍，存在修复空间。
{\small\color{refgray}数据：TFC 2026年净利息收入预期从温和增长下调至持平/微降\par}
{\small\color{refgray}数据：手续费收入增速从高个位数上调至10\%\par}
{\small\color{refgray}数据：2026年ROTCE预期从14\%上调至14\%以上\par}
{\small\color{refgray}数据：长期ROTCE目标16-18\%\par}
{\small\color{refgray}边际变化：市场此前关注净利息收入指引下调，但GS指出背后是资产结构主动优化，非基本面恶化。\par}
\paragraph{GS} 澳大利亚金矿Genesis Minerals与Vault Minerals合并的核心增量来自避免新建Tower Hill选厂，节省约3.5亿澳元资本支出，并规避西澳项目执行风险。合并后近两年黄金产量可提升至约69-78万盎司，全维持成本降至约2850澳元/盎司，较合并前低约10\%；EBITDA预计提升15-25\%。运营现金流回报率从5-10\%提升至10-15\%。
{\small\color{refgray}数据：避免新建选厂节省约3.5亿澳元资本支出\par}
{\small\color{refgray}数据：合并后黄金产量约69-78万盎司/年\par}
{\small\color{refgray}数据：全维持成本降至约2850澳元/盎司，较合并前低约10\%\par}
{\small\color{refgray}数据：EBITDA预计提升15-25\%\par}
{\small\color{refgray}边际变化：市场通常关注合并带来的规模效应，但GS强调价值来源是资本配置路径转变——利用存量产能而非新建。\par}
\paragraph{GS} 中国啤酒行业二季度表现分化：燕京U8 2025年销量约90万千升，2026年一季度仍维持近30\%同比增长，7月后加速。8-10元价格带是高端化核心战场，U8持续从青岛经典、雪花清爽等竞品夺取份额。华润啤酒同期销量低个位数增长，青岛啤酒和重庆啤酒下滑约4\%，百威中国降幅约10\%。
{\small\color{refgray}数据：燕京U8 2025年销量约90万千升\par}
{\small\color{refgray}数据：2026年一季度U8同比增长近30\%\par}
{\small\color{refgray}数据：U8自2019年上市以来年复合增长率约55\%\par}
{\small\color{refgray}数据：华润啤酒二季度销量低个位数增长\par}
{\small\color{refgray}边际变化：市场关注行业总量放缓，但GS指出大单品扩张逻辑未失效，集中资源打透一个价格带更有效。\par}
\paragraph{GS} 中国二季度消费数据出现收入改善与支出观望的背离：居民可支配收入名义同比增速从4.9\%回升至5.6\%，但社零增速从2.4\%降至0.2\%。GS工资追踪指标显示城镇居民工资同比增速从4.3\%升至4.7\%，但居民储蓄率从32.3\%微升至32.6\%，超额存款达61万亿元。消费结构上，食品支出放缓，教育文化娱乐消费回升；以旧换新政策对家电拉动仍在，但对汽车提振减弱。
{\small\color{refgray}数据：居民可支配收入名义同比增速Q1 4.9\% → Q2 5.6\%\par}
{\small\color{refgray}数据：社零同比增速Q1 2.4\% → Q2 0.2\%\par}
{\small\color{refgray}数据：城镇居民工资同比增速Q1 4.3\% → Q2 4.7\%\par}
{\small\color{refgray}数据：居民储蓄率从32.3\%微升至32.6\%\par}
{\small\color{refgray}边际变化：市场关注社零大幅放缓，但GS指出住户调查消费支出增速仅从3.6\%微降至3.8\%，结构变化而非全面恶化。\par}
\paragraph{GS} 贵州茅台2026年7月完成年内第二次调价，出厂价从1269元升至1369元（+7.9\%），i茅台零售价从1539元升至1639元（+6.5\%），1L装飞天零售价上调4.8\%至3269元。两次调价累计出厂价上调17\%，零售价上调9\%，为历史首次一年内两轮调价。GS估算2026年营收增量约1.2\%，利润增量约1.6\%；2027年贡献扩大至约3\%营收和4\%以上利润。调价后原箱飞天批发价从1650元升至1720元。
{\small\color{refgray}数据：出厂价从1269元升至1369元（+7.9\%）\par}
{\small\color{refgray}数据：i茅台零售价从1539元升至1639元（+6.5\%）\par}
{\small\color{refgray}数据：1L装飞天零售价上调4.8\%至3269元\par}
{\small\color{refgray}数据：两次调价累计出厂价+17\%，零售价+9\%\par}
{\small\color{refgray}边际变化：茅台打破过去数年一次调价的节奏，在行业库存调整未完成时主动提价，定价权从市场供需转向公司战略主导。\par}
\paragraph{NOM} 安踏体育二季度运营数据显示，安踏核心品牌和FILA均录得低个位数同比增长，其他品牌（迪桑特、可隆等）增长25-30\%，跑赢运动服饰板块。上半年安踏核心品牌和FILA均实现中个位数增长，其他品牌增速35-40\%。公司已进行系列运营调整：更换安踏核心品牌CEO、推出安踏SV/超级安踏新店型、开设安踏Market大店、狼爪品牌铺开新店。调整后主品牌线上销售恢复双位数增长。
{\small\color{refgray}数据：安踏核心品牌二季度低个位数增长\par}
{\small\color{refgray}数据：FILA二季度低个位数增长\par}
{\small\color{refgray}数据：其他品牌二季度增长25-30\%\par}
{\small\color{refgray}数据：上半年其他品牌增速35-40\%\par}
{\small\color{refgray}边际变化：市场关注行业整体增速放缓，但NOM指出安踏通过主动运营调整（换CEO、新店型）展现出超越平均水平的韧性。\par}
\paragraph{NOM} 老铺黄金二季度销售受金价波动影响，公司已采取主动措施：推出更具吸引力的优惠和赠品、定价更亲民的新品、对高端商场高层级客户提供定向支持，效果初步满意并计划下半年推广。同时持续推进门店升级：深圳万象天地、上海国金中心、上海恒隆广场等门店搬迁或扩建。NOM调整FY26-28F收入和盈利预测5-6\%和9-11\%，反映金价不确定性。
{\small\color{refgray}数据：FY26-28F收入预测下调5-6\%\par}
{\small\color{refgray}数据：FY26-28F盈利预测下调9-11\%\par}
{\small\color{refgray}数据：深圳万象天地、上海国金中心、上海恒隆广场门店升级或扩建\par}
{\small\color{refgray}数据：推出定价更亲民的新品和定向高端客户支持\par}
{\small\color{refgray}边际变化：市场关注金价波动对销售的短期冲击，但NOM强调公司主动应对措施和门店升级投入，认为金价稳定后情绪可能改善。\par}
\subsubsection*{关键数据}
\begin{itemize}
  \item 中国居民可支配收入名义同比增速Q2 5.6\%，社零增速仅0.2\%
  \item 居民储蓄率32.6\%，超额存款61万亿元
  \item 燕京U8 2025年销量90万千升，2026年一季度同比增近30\%
  \item 茅台年内两次调价，出厂价累计+17\%，零售价累计+9\%
  \item TFC 2026年ROTCE预期从14\%上调至14\%以上
  \item 澳洲金矿合并避免新建选厂节省3.5亿澳元，AISC降10\%
  \item 安踏其他品牌二季度增长25-30\%，主品牌线上恢复双位数增长
  \item 老铺黄金FY26-28F盈利预测下调9-11\%
\end{itemize}
\begin{figure}[htbp]
\centering
\includegraphics[width=0.88\linewidth]{figures/fig_042.jpg}
\caption{Exhibit 1 - Exhibit 1: Premium beer SKU map Exhibit 2: We expect premiumization to remain intact in 2025-28E Exhibit 3: Yanjing significantly gained share in the premium segment over the past five years on U8 success}
\end{figure}
\begin{figure}[htbp]
\centering
\includegraphics[width=0.88\linewidth]{figures/fig_045.jpg}
\caption{Exhibit 1 - Exhibit 1: Year-over-year growth in household disposable income per capita rose in Q2 2026 Exhibit 2: Year-over-year growth in household consumption per capita increased slightly in Q2 2026 Exhibit 3: The boost from the consume}
\end{figure}
\begin{figure}[htbp]
\centering
\includegraphics[width=0.88\linewidth]{figures/fig_046.jpg}
\caption{Exhibit 1 - Exhibit 4: Auto sales volume was below last year's level in Q2 2026}
\end{figure}
\begin{figure}[htbp]
\centering
\includegraphics[width=0.88\linewidth]{figures/fig_047.jpg}
\caption{Exhibit 2 - Exhibit 5: Retail sales growth decelerated in Q2 mainly on weaker goods sales}
\end{figure}
{\small\color{refgray}本节综合 7 篇研究；涉及 GS、NOM。}
\newpage
\section{辅助信号｜战略咨询}
本板块聚焦波士顿咨询两份报告：AI嵌入深度决定结构性成本优势，以及国防科技利润结构分析。
\subsection{AI嵌入深度决定结构性成本优势}
\textbf{只有5\%的公司规模化通过AI创造价值，领先者成本降幅是同行三倍，EBIT利润率高出1.6倍，投入资本回报率高出2.7倍。}
\begin{itemize}
  \item AI应用分四阶段：聊天机器人、AI驱动工作流、自主代理、多代理系统；多数公司停留前两阶段，仅捕获小部分价值。
  \item 成本优势在第三、四阶段才显现，AI从执行步骤转向运行决策，价值创造质变。
  \item 领先者与追赶者差距结构化：先发者节省成本持续投入下一轮能力建设，抬高竞争门槛。
  \item 案例：全球消费品公司通过AI部署采购到付款全流程，100天内完成3500次供应商谈判，实现9亿美元成本削减。
  \item 案例：全球制药公司AI应用于小分子药物发现，筛选化合物库扩大100倍，发现周期缩短40\%-50\%。
\end{itemize}
\begin{figure}[htbp]
\centering
\includegraphics[width=0.88\linewidth]{figures/fig_121.jpg}
\caption{EXHIBIT 2 - The durability of profits from sustainment for traditional primes is significant. Defense spending sees large swings based on wartime demand. In the US, for example, defense budgets fell roughly 37\% after the Cold War and 30\% to 45\% in the drawdowns following }
\end{figure}
\subsection{国防科技利润结构：增长快不等于利润大}
\textbf{新赛道增速是传统平台10倍以上，但精尖系统到2033年仍占市场80\%以上份额，利润可持续性远优于新玩家。}
\begin{itemize}
  \item 国防航空技术分三类：精尖系统（如F/A-18）、可负担量产系统（如协作作战飞机）、可消耗系统（如巡飞弹）。
  \item 2025年精尖系统采购收入约650亿美元，可负担量产约50亿美元，可消耗仅5500万美元；到2033年精尖系统年增长2\%-3\%，可消耗增长35\%-40\%。
  \item 精尖系统一半利润来自后期维护（如F/A-18E/F全生命周期利润超10亿美元），新赛道利润池小且一次性。
  \item 传统平台利润可持续性源于数十年服务合同，新玩家利润受研发不确定性影响大。
\end{itemize}
\begin{figure}[htbp]
\centering
\includegraphics[width=0.88\linewidth]{figures/fig_121.jpg}
\caption{EXHIBIT 2 - The durability of profits from sustainment for traditional primes is significant. Defense spending sees large swings based on wartime demand. In the US, for example, defense budgets fell roughly 37\% after the Cold War and 30\% to 45\% in the drawdowns following }
\end{figure}
\newpage
\section{辅助信号｜智库/国际机构}
本板块整合BIS、IMF、OECD、世界银行最新研究，聚焦三大主题：AI重塑网络风险与公共治理、全球供应链与产业政策的结构性瓶颈、以及制度质量对宏观与发展的深层影响。
\subsection{AI重塑网络风险与公共治理}
\textbf{前沿AI模型（如Mythos、GPT-5.5）正成为网络攻防的转折点，攻击成本快速下降且能力跃升；同时，公共部门在AI应用上存在“指南多、操作少”的断层，影子AI蔓延，结构化实验成为必要路径。}
\begin{itemize}
  \item BIS指出，Mythos在专家级网络攻防任务中通过率达68.6\%，GPT-5.5达71.4\%；一次完整攻击链成本低至50-100美元（使用DeepSeek），攻击方成本优势显著。
  \item Mythos发布后，Firefox月度安全漏洞修复量跃升六倍，每日新增关键级CVE（CVSS≥9.0）从2022-2025年的10个升至2026年4月后的近20个，防御压力陡增。
  \item BIS建议，系统性重要资金基础设施应参与新AI模型预测试，并重新校准跨境信息共享机制，以应对AI驱动的攻击速度与规模。
  \item OECD报告显示，14国中仅少数（如英国、澳大利亚）提供详细AI实验操作指南，22\%的公共部门受访者已使用生成式AI，但大量为未经审批的“影子AI”。
  \item OECD强调，结构化实验是政府在概率性技术面前建立能力、管理不确定性并赢得公众信任的必要步骤，需包括试点设计、评估指标和失败实验处理机制。
\end{itemize}
\begin{figure}[htbp]
\centering
\includegraphics[width=0.88\linewidth]{figures/fig_109.jpg}
\caption{Figure 5 - The analysis points to highly uneven distribution of policy attention across different dimensions of AI development and implementation in the public sector. While data governance, AI design principles, public-private partnerships, and investments in skills, li}
\end{figure}
\subsection{全球供应链与产业政策的结构性瓶颈}
\textbf{稀土供应链瓶颈在冶炼分离环节（中国占比超85\%），而非采矿；AI算力出口优势不来自廉价电力，而受制于硬件统一定价与制度摩擦；产业政策需精准瞄准瓶颈，否则成本高而效果有限。}
\begin{itemize}
  \item IMF量化显示，全球稀土冶炼分离环节中国产能占比超85\%，非中国地区采矿潜力虽大但冶炼能力建设需更长周期和更高资本支出。
  \item 模拟表明，通过CAPEX补贴实现适度去风险（如自给率从10\%提至25\%），美国单干十年总成本约11.3亿美元，联合日欧澳可降至9亿美元；但去风险程度越高，边际成本上升越快。
  \item 世界银行模型揭示，算力单位成本中硬件摊销占约90\%，全球统一定价导致各国成本差异仅12-20\%；融资成本差距（WACC每差10个百分点，单位成本增0.29美元）可抹平电价优势。
  \item 双边制度摩擦（数据治理可信度、监管稳定性、地缘政治信任）足以消除廉价电力带来的成本优势，马来西亚、印度、巴西等吸引最多投资的发展中地区在成本排名中仅处中游。
  \item IMF指出，数字服务税（DST）与跨境服务贸易下降存在统计关联，税负可能转嫁给消费者；目的地型增值税被列为“最连贯、扭曲最小”的方案，但实施需全球协调以避免税制重叠。
\end{itemize}
\begin{figure}[htbp]
\centering
\includegraphics[width=0.88\linewidth]{figures/fig_089.jpg}
\caption{Figure 1 - Figure 1: Rare Earth Mining by Country 1998–2024 (a) LREE Production (b) HREE Production Sources: Nassar et al. (2023); U.S. Geological Survey (2026); and IMF staff calculations. Note: Production is measured in thousands of me}
\end{figure}
\begin{figure}[htbp]
\centering
\includegraphics[width=0.88\linewidth]{figures/fig_115.jpg}
\caption{Figure 1 - WonderNetwork. (2024). Global Ping Statistics. wondernetwork.com. World Bank. (2024). World Development Indicators. Washington, DC. World Bank. (2025). Digital Progress and Trends Report 2025: Strengthening AI Foundations. Washington, DC: World Bank. Figure 1.}
\end{figure}
\subsection{制度质量对宏观与发展的深层影响}
\textbf{汇率制度、反洗钱框架、正规化政策等制度因素，通过收入分配、资本流动和社会资本积累等渠道，对宏观经济稳定与发展产生系统性影响，其效果常被传统指标掩盖。}
\begin{itemize}
  \item IMF对吉尔吉斯斯坦的研究显示，结构性流动性过剩（占GDP 9.2\%）使市场利率长期低于政策利率，走廊下限成为实际锚，单一盯住政策利率可能高估紧缩力度。
  \item IMF反洗钱指引首次将AML/CFT条件性纳入贷款项目设计，认为洗钱可通过资本外流、银行渗透和国际收支数据失真对宏观产生冲击，未来反洗钱改革或成贷款硬约束。
  \item 世界银行基于哥伦比亚移民合法化项目的中期证据表明，合法身份带来的社会资本收益（融入指数高0.95个标准差）需四年半才显现，且依赖“稳定身份+日常正面互动”。
  \item 世界银行前行长Malpass批评当前浮动汇率与多重汇率体系，以埃塞俄比亚为例：2019年比尔兑美元28:1，2025年跌至161:1，贫困率从33\%升至43\%，贬值成为对低收入群体的隐性征税。
  \item IMF对葡萄牙的研究显示，卫生支出增长主因是“更贵”而非“更多”，人员费用和药品价格是主要驱动；生产率差距的根源在企业融资结构，欧洲企业研发投入仅为美国同行的五分之三。
\end{itemize}
\begin{figure}[htbp]
\centering
\includegraphics[width=0.88\linewidth]{figures/fig_077.jpg}
\caption{Figure 1 - This study assesses the effective monetary stance in the Kyrgyz Republic amid persistent excess liquidity, rapid household credit growth, and elevated inflation pressures. It constructs a monthly Financial Conditions Index (FCI) that captures liquidity conditi}
\end{figure}
\begin{figure}[htbp]
\centering
\includegraphics[width=0.88\linewidth]{figures/fig_116.jpg}
\caption{Figure 1 - Figure 1: First-Stage Discontinuity in PEP-RAMV Take-Up at the Program-Eligibility Cut-off Notes: Panel (a) plots the weekly mean of the household-level PEP-RAMV take-up indicator (navy line, left axis) against date of arrival in}
\end{figure}
\subsection{AI、技术与生产率}
\textbf{用于判断 AI 投资周期、生产率扩散和产业约束是否正在改变资产定价。}
\begin{itemize}
  \item 衡量农业是否在环境可持续的前提下实现增长，长期以来缺乏一个公认的标尺。传统全要素生产率只算产出和投入，把温室气体、氨排放、氮磷流失等环境外部性排除在外。经合组织最新发布的这份工作论文，系统比较了六种“环境可持续生产率增长”指标，并用墨西哥、荷兰、新西兰三个国家的三十年数据做了实证...
\end{itemize}
\begin{figure}[htbp]
\centering
\includegraphics[width=0.88\linewidth]{figures/fig_110.jpg}
\caption{Figure 2 - Figure 2 (panel b). Between 2000 and the global financial crisis, there was a substantial increase in international fertiliser prices, with urea prices increasing fivefold between 1999 and 2007 and DAP prices more than doubling (World Bank, 2025[55]). Increasi}
\end{figure}
\subsection{宏观政策与金融稳定}
\textbf{用于校准利率、通胀、财政约束和金融稳定叙事。}
\begin{itemize}
  \item 吉尔吉斯斯坦约65\%的劳动力在非正式部门就业，非正式宏观环境贡献了约19\%的GDP。IMF最新发布的国别报告指出，这个中亚国家的非正式宏观环境并非简单的“逃税”问题，而是一个由教育质量、制度碎片化和社保缺位共同塑造的结构性困境。报告的核心判断是：降低非正式化，不能靠强制取缔，而应...
\end{itemize}
\begin{figure}[htbp]
\centering
\includegraphics[width=0.88\linewidth]{figures/fig_094.jpg}
\caption{Figure 1 - Figure 2: Population Tree of the Kyrgyz Republic in 2024.}
\end{figure}
\section{结语}
后续需验证的数据与叙事节点包括：美国核心PCE实际走势与市场预期的偏差、中东局势对能源价格的持续性影响、中国三季度财政加速的力度与效果、AI资本开支回报率的初步数据、以及中国电池消费税对二线厂商出清节奏的实际影响。
\newpage
\section{覆盖概览}
投行/券商 31 篇；战略咨询 2 篇；智库/国际机构 13 篇。\par
投行覆盖：GS 17篇 / Bernstein 6篇 / NOM 5篇 / JPM 2篇 / MS 1篇\par
\section{Disclaimer}
{\scriptsize\color{refgray}
This community is a paid learning and information-sharing community created for general educational purposes only. The membership fee is charged solely for access to community discussions, educational content, organization, and operational costs. It is not an advisory fee, management fee, performance fee, brokerage fee, or compensation for personalized financial, investment, legal, tax, or professional advice.\par
I participate and share information solely as an individual and for educational discussion among community members. I am not acting as a financial adviser, investment adviser, broker, dealer, portfolio manager, legal adviser, tax adviser, accountant, fiduciary, or any other licensed professional adviser. Nothing shared in this community constitutes, and should not be interpreted as, financial advice, investment advice, legal advice, tax advice, a recommendation, solicitation, offer, endorsement, instruction, or guarantee to buy, sell, hold, short, trade, or invest in any security, cryptocurrency, fund, derivative, strategy, or other financial product.\par
All content shared in this community, including messages, discussions, links, files, charts, examples, market commentary, watchlists, AI-generated or AI-polished summaries, and third-party materials, is general, impersonal, and educational in nature. It does not take into account any individual's financial situation, investment objectives, risk tolerance, investment horizon, liquidity needs, tax status, legal restrictions, or suitability requirements. No content should be relied upon as suitable for any specific person, account, portfolio, or financial decision.\par
Information shared in this community may be incomplete, inaccurate, outdated, speculative, AI-generated, AI-edited, or based on third-party internet sources that have not been independently verified. I make no representation or warranty as to the accuracy, completeness, timeliness, reliability, or suitability of any information shared.\par
Financial markets involve substantial risk, including the possible loss of principal. Past performance, historical data, hypothetical examples, simulated results, backtests, personal opinions, or educational examples do not guarantee future results. Each member is solely responsible for conducting their own research, exercising independent judgment, and making their own decisions. Before making any financial, investment, legal, or tax decision, members should consult qualified and licensed professionals.\par
Participation in this community, payment of any membership fee, or communication with me or other members does not create any adviser-client, fiduciary, professional, contractual, agency, partnership, employment, or other advisory relationship. By joining or participating in this community, you acknowledge and agree that you are solely responsible for your own decisions, actions, transactions, profits, losses, risks, and consequences, and that I shall not be liable for any direct, indirect, incidental, consequential, financial, legal, tax, or other loss or damage arising from or related to any content, discussion, information, or material shared in this community.\par
}
\newpage
\begin{center}
{\Large 更多详情报告kcdesk.com}\par\vspace{0.8cm}
\includegraphics[width=0.58\linewidth]{\detokenize{../../prompts/zsxq_img.jpg}}
\end{center}
\end{document}
